\documentclass[journal]{IEEEtran}

\usepackage{amsmath}
\usepackage{booktabs}
\usepackage{array}
\usepackage{url}
\usepackage{multirow}
\usepackage{makecell}
\usepackage{hyperref}
\hypersetup{colorlinks=true, linkcolor=blue, citecolor=blue, urlcolor=blue}

\hyphenation{op-tical net-works semi-conduc-tor re-si-li-ence}

\begin{document}

%% ─── TITLE ────────────────────────────────────────────────────────────────────
\title{CRASP: A Machine Learning-Driven Cyber Resilience Assessment and
Scoring Platform Integrating 2026 Threat Intelligence}

\author{Jenah~Karthick~Palanikumar%
\thanks{Jenah Karthick Palanikumar is with the School of Computer Science,
University of Galway, Galway, Ireland
(e-mail: j.palanikumar1@universityofgalway.ie).}%
\thanks{Manuscript submitted August 2026.}}

\markboth{MSc Computer Science Capstone Project,~University of Galway,~August~2026}%
{Palanikumar: CRASP --- Cyber Resilience Assessment and Scoring Platform}

\maketitle

%% ─── ABSTRACT ─────────────────────────────────────────────────────────────────
\begin{abstract}
The growing complexity and cost of cyber threats---with the average breach
now costing \$4.99M globally (IBM, 2026)---demand quantitative,
evidence-based tools for organisational resilience assessment. This paper
presents CRASP (Cyber Resilience Assessment and Scoring Platform), a
machine learning-driven framework that evaluates cyber resilience across five
dimensions derived from NIST~CSF~2.0: Anticipate, Withstand, Recover, Adapt,
and Evolve. CRASP addresses critical limitations of prior work through:
(1)~a \emph{latent variable design} that prevents data leakage and produces
honest predictive uncertainty (theoretical $R^2$ ceiling~$\approx 0.87$);
(2)~integration of 2026 threat intelligence from Verizon DBIR~2026, IBM Cost
of a Data Breach 2026, MITRE ATT\&CK~v19, NIST~CSF~2.0, and Gartner~2026;
(3)~SMOTE-based class balancing resolving Critical-class recall from 0\%
to~38\%; (4)~SHAP-based explainability for per-prediction feature attribution;
and (5)~a What-If intervention simulator enabling evidence-based ROI analysis
across 12 security interventions. Trained on 3,000 synthetic organisations
with 80+ features, CRASP achieves 80\% classification accuracy (Macro
AUC~$= 0.942$), average regression $R^2 = 0.869$, NLP CVE severity accuracy
of 91\%, and anomaly detection with a $4\times$ breach-rate gap between
anomalous and normal organisations. An interactive Streamlit dashboard
provides real-time assessment, SHAP waterfall explanations, trajectory
forecasting, PDF reporting, and scenario planning. Results demonstrate that
combining 2026 threat intelligence with explainable ML produces a
defensible, practically valuable, and academically honest cyber resilience
assessment instrument.
\end{abstract}

\begin{IEEEkeywords}
Cyber resilience, machine learning, XGBoost, SHAP explainability, SMOTE,
NIST CSF 2.0, threat intelligence, anomaly detection, natural language
processing, Streamlit, latent variable design.
\end{IEEEkeywords}

\IEEEpeerreviewmaketitle


%% ─── 1. INTRODUCTION ──────────────────────────────────────────────────────────
\section{Introduction}

\IEEEPARstart{T}{he} cyber threat landscape has undergone a fundamental
transformation that existing assessment tools are not equipped to address.
The Verizon Data Breach Investigations Report 2026 (DBIR 2026)~\cite{verizon2026}
identifies supply chain breaches as the fastest-growing attack vector,
now accounting for 48\% of all incidents---a 60\% year-over-year increase.
The median time to patch critical vulnerabilities has risen to 43~days
despite widely available patch management tooling, and shadow artificial
intelligence (AI) exposure now affects an estimated 45\% of the enterprise
workforce, creating an unmonitored attack surface that did not exist in prior
generation security models. Simultaneously, the IBM Cost of a Data Breach
Report 2026~\cite{ibm2026} places the global average breach cost at \$4.99
million---a 12\% increase over the prior year---with Healthcare organisations
sustaining \$6.64 million, Financial Services \$6.08 million, and Technology
\$5.62 million. Organisations in sectors that adopt AI security tools reduced
their breach cost by an average of \$2.2 million, yet adoption remains at
only 34\% of large enterprises~\cite{ibm2026}.

In this environment, organisations require assessment tools that go beyond
annual compliance audits and checklist-based frameworks. They need
\emph{quantitative, predictive, and continuously updated} cyber resilience
intelligence that can translate measurable security controls into a defensible
score, explain why that score was assigned, and simulate how specific
interventions would change it.

Existing frameworks such as NIST~CSF~2.0~\cite{nist_csf2},
ISO/IEC~27001:2022~\cite{iso27001}, and MITRE ATT\&CK~v19~\cite{mitre_attack}
provide valuable structural guidance but lack quantitative scoring, predictive
capability, and integration with 2026 threat data. Commercial tools---including
BitSight, SecurityScorecard, and RiskRecon---offer scoring but are proprietary,
opaque in methodology, and not calibrated against 2026 industry statistics.
Academic implementations frequently suffer from circular label-generation
methodologies that inflate reported performance metrics, as documented in
Section~\ref{sec:related}.

This paper presents CRASP (Cyber Resilience Assessment and Scoring Platform),
which addresses these gaps through the following principal contributions:
\begin{enumerate}
  \item A five-dimension resilience scoring model (Anticipate, Withstand,
        Recover, Adapt, Evolve) grounded in NIST~CSF~2.0 and calibrated
        against 2026 industry statistics from five authoritative reports;
  \item A \emph{latent variable design} ensuring academic integrity by
        withholding an unobservable cultural score ($S_\text{culture}$) from
        all ML models---a methodological flaw identified and corrected from
        a prior circular implementation;
  \item Integration of six concrete features derived from 2026 threat
        intelligence: shadow AI exposure, supply chain vendor count, AI
        security tool adoption, average patch days, AI-enabled breach history,
        and estimated sector-adjusted breach cost;
  \item SHAP~\cite{lundberg2017shap} post-hoc explainability transforming
        black-box predictions into per-prediction feature attribution waterfall
        charts rendered in a real-time dashboard;
  \item A What-If intervention simulator evaluating 12 security interventions
        with cost, timeline, and ROI estimates; and
  \item Rigorous multi-model empirical evaluation including paired $t$-tests,
        multiclass ROC/AUC analysis, learning curves, calibration analysis,
        and sector-level fairness measurement.
\end{enumerate}

The remainder of this paper is structured as follows. Section~\ref{sec:related}
reviews related work and identifies gaps. Section~\ref{sec:method} presents
the system methodology. Section~\ref{sec:results} reports empirical results.
Section~\ref{sec:disc} discusses findings, implications, and limitations.
Section~\ref{sec:conc} concludes with future directions.


%% ─── 2. RELATED WORK ──────────────────────────────────────────────────────────
\section{Literature Review}
\label{sec:related}

\subsection{Cyber Resilience Frameworks}

The NIST Cybersecurity Framework 2.0~\cite{nist_csf2} introduced a revised
six-function structure (Govern, Identify, Protect, Detect, Respond, Recover)
with the notable addition of the \emph{Govern} function, reflecting recognition
that cybersecurity risk management must be embedded in organisational governance
rather than treated as a purely technical problem. While CSF~2.0 provides
a comprehensive vocabulary and implementation tiers, it prescribes no
quantitative scoring methodology. The associated AI Risk Management
Framework~\cite{nist_ai_rmf} extends CSF~2.0 principles to AI-specific risks,
directly informing CRASP's treatment of shadow AI exposure and AI-enabled
breach history as quantitative features.

ISO/IEC~27001:2022~\cite{iso27001} provides the most widely adopted ISMS
standard globally, with Annex~A mapping 93 controls across four themes
(Organisational, People, Physical, Technological). These controls correspond
broadly to CRASP's Withstand (preventive controls) and Evolve (improvement
and learning) dimensions. However, ISO~27001 certification produces a binary
pass/fail outcome rather than a continuous resilience score, and re-certification
cycles of two to three years mean the assessment does not reflect the
organisation's current posture.

MITRE ATT\&CK~v19~\cite{mitre_attack} provides the most granular taxonomy
of adversary techniques (858 enterprise techniques across 14 tactics), which
CRASP uses for CVE contextualisation and recommendation alignment. The D3FEND
knowledge graph~\cite{nist_csf2} provides complementary defensive technique
mappings that inform CRASP's recommendation engine, enabling output such as
``apply network segmentation to mitigate lateral movement via T1021 (Remote
Services).''

\subsection{Machine Learning for Cyber Risk Quantification}

XGBoost~\cite{chen2016xgboost}, the gradient boosted tree ensemble introduced
by Chen and Guestrin~(2016), has consistently demonstrated state-of-the-art
performance on structured tabular classification benchmarks, including the
majority of winning entries in Kaggle competitions involving tabular data.
Its key advantages---regularised tree growth, built-in handling of missing
values, and column subsampling---make it well-suited to the heterogeneous
feature space of cyber resilience assessment, which combines continuous metrics
(e.g., \texttt{security\_budget\_pct}), binary indicators (e.g.,
\texttt{uses\_siem}), and categorical variables (e.g., sector, size).

SHAP (SHapley Additive exPlanations)~\cite{lundberg2017shap} provides the
most theoretically grounded approach to post-hoc ML explainability, satisfying
three critical axioms: \emph{local accuracy} (SHAP values sum to the
prediction), \emph{missingness} (absent features have zero attribution), and
\emph{consistency} (a feature cannot have lower attribution in a model that
uses it more). SHAP TreeExplainer computes exact SHAP values for tree-based
models in polynomial rather than exponential time, making it practical for
the 80-feature CRASP model.

SMOTE~\cite{chawla2002smote} addresses the fundamental class imbalance problem
that arises in any real cyber resilience dataset: organisations with
\emph{Critical} risk posture are rare (by definition, they are extreme
outliers), making them severely under-represented in any unbiased sample.
SMOTE generates synthetic minority-class samples through linear interpolation
between existing minority samples in feature space, rather than simple
oversampling, which would cause overfitting.

\subsection{Prior Academic Implementations}

Several prior works apply ML to cyber risk quantification with relevant
findings for CRASP's design. Jakubik et~al.~\cite{jakubik2021} applied
random forests to breach prediction using 47 financial and operational
features, reporting AUC of 0.82 but providing no explainability mechanism
and not addressing 2026 threat dimensions. Fielder et~al.~\cite{fielder2016}
proposed a game-theoretic approach to cyber investment optimisation using
multi-objective optimisation over a simplified control set, achieving
theoretical Pareto-optimal security portfolios but lacking real-time scoring
and empirical validation at scale. Srinivas et~al.~\cite{srinivas2019}
applied neural networks to vulnerability assessment, reporting $R^2 = 0.94$,
but their \emph{circular label-generation methodology}---where ground-truth
labels were derived directly from the same observable features used for
training---renders the reported $R^2$ trivially high. Any model that
memorises the label formula would achieve similarly inflated metrics.
CRASP directly corrects this flaw through the latent variable design
described in Section~\ref{sec:latent}.

\subsection{Gaps in Existing Approaches}

Table~\ref{tab:gap} summarises how CRASP compares to the most relevant
existing tools and frameworks across six dimensions. No existing open-source
tool simultaneously addresses all six: quantitative scoring, ML-based
prediction, SHAP explainability, 2026 threat intelligence integration,
intervention simulation, and academically honest label design.

\begin{table}[!t]
\renewcommand{\arraystretch}{1.3}
\caption{Comparison of Cyber Resilience Assessment Approaches}
\label{tab:gap}
\centering
\small
\begin{tabular}{lcccccc}
\toprule
\textbf{Tool / Work} & \textbf{Quant.} & \textbf{ML} & \textbf{Expl.}
  & \textbf{2026} & \textbf{Sim.} & \textbf{Honest} \\
\midrule
NIST CSF 2.0~\cite{nist_csf2}        & \texttimes & \texttimes & \texttimes & \texttimes & \texttimes & \checkmark \\
ISO 27001~\cite{iso27001}             & \texttimes & \texttimes & \texttimes & \texttimes & \texttimes & \checkmark \\
BitSight (commercial)                 & \checkmark & \checkmark & \texttimes & \texttimes & \texttimes & ? \\
Jakubik et al.~\cite{jakubik2021}     & \checkmark & \checkmark & \texttimes & \texttimes & \texttimes & \checkmark \\
Fielder et al.~\cite{fielder2016}     & \checkmark & \texttimes & \texttimes & \texttimes & \checkmark & \checkmark \\
Srinivas et al.~\cite{srinivas2019}   & \checkmark & \checkmark & \texttimes & \texttimes & \texttimes & \texttimes \\
\textbf{CRASP (this work)}            & \checkmark & \checkmark & \checkmark & \checkmark & \checkmark & \checkmark \\
\bottomrule
\multicolumn{7}{l}{\footnotesize Quant.=Quantitative; Expl.=Explainable; 2026=2026 Intel.; Sim.=Simulator}
\end{tabular}
\end{table}


%% ─── 3. METHODOLOGY ───────────────────────────────────────────────────────────
\section{Methodology}
\label{sec:method}

\subsection{System Architecture}

CRASP follows a four-layer pipeline architecture executed by a single
\texttt{run.sh} orchestration script:

\textbf{Layer~1 --- Data:} The \texttt{generate\_data.py} module produces
3,000 synthetic organisational profiles, 1,500 CVE records, and 300
organisations with 18-month time-series histories. Data generation is seeded
for reproducibility and parameterised from 2026 industry statistics.

\textbf{Layer~2 --- Feature Engineering:} The \texttt{train\_models.py}
module extracts 80+ features per organisation, applies
\texttt{StandardScaler} normalisation for regression and NLP tasks, and
computes derived features including a normalised patch urgency score
(\texttt{avg\_patch\_days} $\times$ \texttt{critical\_cve\_count} / 43) and
an AI risk exposure index combining \texttt{shadow\_ai\_exposure} and
\texttt{ai\_breach\_history}.

\textbf{Layer~3 --- ML Models:} Five distinct models are trained and
serialised to \texttt{data/models/} using \texttt{joblib}. All metrics are
written to \texttt{model\_metrics.json} for consumption by the dashboard.

\textbf{Layer~4 --- Dashboard:} A 7-tab Streamlit application
(\texttt{dashboard/app.py}) provides interactive assessment, SHAP
visualisations, trend forecasting, PDF export via fpdf2, and What-If
scenario simulation. All paths use Python's \texttt{pathlib.Path} for
cross-platform compatibility.

\subsection{Data Generation and Latent Variable Design}
\label{sec:latent}

A corpus of 3,000 synthetic organisations was generated across eight industry
sectors and four size categories. Table~\ref{tab:features} details the six
2026-specific features added in this work, alongside their source statistics
and role in the scoring formula. Standard features include MFA coverage,
patch compliance, security budget percentage, incident response time (hours),
penetration test frequency, security staff FTE, training completion rate,
SIEM/SOAR/EDR deployment flags, and prior breach count---80 features in total.

\begin{table}[!t]
\renewcommand{\arraystretch}{1.3}
\caption{2026-Specific Features: Source Statistics and Role}
\label{tab:features}
\centering
\small
\begin{tabular}{p{2.1cm}p{2.2cm}p{3.0cm}}
\toprule
\textbf{Feature} & \textbf{2026 Source} & \textbf{Role in CRASP} \\
\midrule
\texttt{avg\_patch\_days} & DBIR 2026: median 43~days & Anticipate penalty; NLP CVE context \\
\texttt{shadow\_ai\_exposure} & DBIR 2026: 45\% workforce & Evolve risk factor; reduces score \\
\texttt{supply\_chain\_vendors} & DBIR 2026: 48\% vector & Withstand exposure index \\
\texttt{uses\_ai\_security\_tools} & IBM 2026: saves \$2.2M & Detect/respond capability flag \\
\texttt{ai\_breach\_history} & IBM 2026: 25\% of breaches & Classification label shift \\
\texttt{estimated\_breach\_cost} & IBM 2026: \$4.99M avg & Risk quantification output \\
\bottomrule
\end{tabular}
\end{table}

The central methodological contribution is the \textbf{latent variable
design}. Each organisation's true resilience score is determined by:
\begin{equation}
  S_{\text{true}} = 0.78\,S_{\text{formula}} + 0.22\,S_{\text{culture}}
                    + \varepsilon, \quad \varepsilon \sim \mathcal{N}(0,\,4)
  \label{eq:latent}
\end{equation}
where $S_{\text{formula}}$ is computed deterministically from observable
features using sector-weighted control scores, and $S_{\text{culture}}$ is an
\emph{unobservable} latent variable representing security leadership quality,
team morale, and informal knowledge sharing---factors that are acknowledged
in industry research~\cite{verizon2026} as significant determinants of
resilience outcomes but that cannot be directly measured from a security
control inventory. $S_{\text{culture}}$ is sampled from
$\mathcal{N}(\mu_\text{sector}, \sigma_\text{sector})$ where parameters are
calibrated per sector (e.g., Finance: $\mu=72$, $\sigma=8$; Retail:
$\mu=54$, $\sigma=12$).

Critically, $S_{\text{culture}}$ is generated at data-creation time but is
\emph{never provided to any ML model as a training feature}. This ensures
that models face genuine predictive uncertainty---the 22\% latent contribution
to $S_{\text{true}}$ is irreducibly unobservable from the feature set---placing
a theoretical $R^2$ ceiling at $\approx 0.87$. The prior implementation (v1)
omitted this design and achieved a spurious $R^2 = 0.93$ by allowing models
to trivially reproduce the scoring formula from its inputs, a circular pattern
that has also been identified in Srinivas et~al.~\cite{srinivas2019}.

\subsection{ML Model Suite}

\subsubsection{Classification}
XGBoost~\cite{chen2016xgboost} classifies organisations into five resilience
levels: \emph{Critical} ($<$40), \emph{Developing} (40--54),
\emph{Managed} (55--69), \emph{Advanced} (70--84), and \emph{Optimized}
($\geq$85). The class distribution in the raw dataset reflects the latent
variable design: Critical (1.4\%), Developing (18.2\%), Managed (47.6\%),
Advanced (28.4\%), Optimized (4.4\%). SMOTE~\cite{chawla2002smote} was
applied exclusively to the training set (80\% split) to produce 6,725
balanced examples across all five classes from the original 2,400, ensuring
no synthetic test contamination. XGBoost hyperparameters were tuned via
randomised search: \texttt{n\_estimators=400}, \texttt{max\_depth=5},
\texttt{learning\_rate=0.05}, \texttt{subsample=0.8},
\texttt{colsample\_bytree=0.7}, \texttt{min\_child\_weight=3},
\texttt{gamma=0.1}, \texttt{reg\_alpha=0.1}, \texttt{reg\_lambda=1.5}.
Per-sample SHAP-weighted \texttt{sample\_weight} was applied to further
emphasise minority class examples. Baselines are Random
Forest~\cite{breiman2001} (\texttt{n\_estimators=300},
\texttt{class\_weight=`balanced'}) and Logistic Regression (multinomial,
$C = 0.5$).

\subsubsection{Regression}
Five independent XGBoost regressors predict each CRASP dimension score
on a 0--100 scale: Anticipate (threat awareness), Withstand (resistance
controls), Recover (response capability), Adapt (learning agility), and
Evolve (strategic improvement). The 22\% latent variance from
Equation~(\ref{eq:latent}) sets the theoretical $R^2$ ceiling at $\approx 0.87$.
Separate Random Forest regressors serve as baselines for each dimension.

\subsubsection{Anomaly Detection}
An Isolation Forest~\cite{liu2008isolation} (\texttt{contamination=0.05},
\texttt{n\_estimators=300}, \texttt{random\_state=42}) identifies
organisations with statistically unusual security profiles in the 80-dimensional
feature space. The Isolation Forest operates by random feature selection and
random split-point selection, isolating anomalous instances with fewer
partitions on average than normal instances. Validation compares breach rates,
mean resilience scores, and security spending patterns between anomalous and
normal groups to confirm practical detector utility.

\subsubsection{Time-Series Forecasting}
A Gradient Boosting Regressor (GBR, \texttt{n\_estimators=200},
\texttt{max\_depth=4}, \texttt{learning\_rate=0.1}) predicts 6- and
12-month resilience trajectories using sliding-window features extracted from
300 organisations observed over 18 monthly time steps. Window features include
a 3-month rolling mean score, 3-month rolling standard deviation (volatility),
linear trend slope, recent improvement/decline flag, and sector and size
dummy variables. Prophet~\cite{taylor2018prophet} was evaluated as an
alternative and yielded $R^2 = -88.6$---substantially below even a
constant-prediction baseline. This confirms that Prophet's Fourier-series
seasonal decomposition requires multi-year histories with clear seasonal
patterns, neither of which characterise this 18-month synthetic dataset.
This negative empirical finding is reported transparently as it informs the
model selection rationale.

\subsubsection{NLP --- CVE Severity Classification}
A TF-IDF + LinearSVC~\cite{sklearn2011} pipeline classifies 1,500 CVE
descriptions into CRITICAL/HIGH/MEDIUM/LOW severity. Labels are derived
from CVSS v3.1 numeric scores (CRITICAL: $\geq$9.0; HIGH: 7.0--8.9;
MEDIUM: 4.0--6.9; LOW: $<$4.0), \emph{not} from keyword matching, ensuring
the classifier must learn genuine semantic patterns rather than memorise
label-correlated keywords. Three severity template tiers create realistic
prediction difficulty: 40\% of descriptions include severity-discriminating
vocabulary (CVSS-aligned qualifiers such as ``unauthenticated,
zero-click remote code execution''); 35\% use partial attack-type context
(e.g., ``SQL injection vulnerability'') without severity-specific framing;
and 25\% use fully generic boilerplate text (``A vulnerability exists that
could allow an attacker to...''), producing a realistic 9\% irreducible
error rate. Competing baselines are TF-IDF + Logistic Regression and
TF-IDF + Random Forest.

\subsection{Explainability, Calibration, and Evaluation Protocol}

Post-hoc explainability uses SHAP TreeExplainer~\cite{lundberg2017shap}
applied to the XGBoost classifier. Global feature importance is computed as
the mean absolute SHAP value $\bar{|\phi_j|}$ across all test-set predictions,
providing a model-wide ranking of feature influence. At the individual
assessment level, the dashboard renders per-prediction SHAP waterfall charts
that decompose any prediction into additive feature contributions, coloured
green (positive impact) or red (negative impact) relative to the base
expected value.

A Calibrated Classifier (CalibratedClassifierCV with isotonic regression,
3-fold) is trained on held-out training data to post-process XGBoost's
probability outputs into reliably calibrated confidence estimates. This
ensures that a stated ``79\% confidence in Advanced level'' corresponds to
the model being correct approximately 79\% of the time across similarly
confident predictions---a property essential for risk communication to
non-technical stakeholders.

Model evaluation follows a rigorous protocol: stratified 80/20 train-test
split with \texttt{random\_state=42} for reproducibility; 5-fold
cross-validation with stratified splitting; one-vs-rest multiclass ROC curves
with per-class AUC; paired $t$-test on 5-fold CV folds for XGBoost vs.\
Random Forest statistical comparison; learning curves at 10 training set
fractions with 95\% confidence intervals via 5-fold bootstrap; and
sector-level fairness analysis computing per-sector classification accuracy
across all eight sectors.

\subsection{Dashboard Design}

The Streamlit dashboard provides seven tabs serving distinct user roles:
(1)~\textbf{Organisation Assessment}: form-based input with SHAP waterfall
and PDF download; (2)~\textbf{Threat Landscape}: live 2026 statistics with
breach cost calculator; (3)~\textbf{Risk Intelligence}: CVE severity
classification and anomaly flagging; (4)~\textbf{Forecasting}: 6- and
12-month trajectory projection; (5)~\textbf{Benchmarking}: sector and size
peer comparison; (6)~\textbf{Model Performance}: ROC curves, learning
curves, SHAP global importance, sector fairness, and model comparison
tables; and (7)~\textbf{What-If Simulator}: 12-intervention scenario
planner with radar chart and ROI estimate.

PDF reports are generated via \texttt{fpdf2}, containing the organisation
profile, overall score, per-dimension goal table, top-8 recommendations
aligned to MITRE ATT\&CK~v19 techniques, and 2026 threat indicators relevant
to the organisation's sector. Reports are downloadable directly from the
browser without server-side storage.

\subsection{What-If Simulator Design}

The What-If simulator enables practitioners to plan security investment
portfolios before committing budget. Twelve interventions are modelled, each
with: (a)~a feature delta applied to the current organisation profile (e.g.,
SIEM deployment sets \texttt{uses\_siem=True} and reduces
\texttt{mean\_time\_to\_detect} by 40\%); (b)~an estimated implementation
cost in USD; and (c)~an estimated deployment timeline in weeks. After
selecting interventions, the simulator applies all feature deltas, runs the
updated profile through all five regression models and the classifier,
and computes the projected new overall score and dimension scores. Results
are displayed as a before/after radar chart across the five dimensions, a
per-dimension bar chart showing score improvements, and an ROI estimate
calculated as ($\Delta$score $\times$ \$100,000) / intervention\_cost---a
normalised metric expressing score improvement per \$100,000 invested.


%% ─── 4. RESULTS ───────────────────────────────────────────────────────────────
\section{Results and Analysis}
\label{sec:results}

\subsection{Classification Performance}

Table~\ref{tab:classif} presents results for the three-model comparison across
four evaluation metrics. XGBoost achieves 80.0\% test accuracy and Macro
AUC~$= 0.942$, outperforming Random Forest and Logistic Regression on test
accuracy while all three models show comparable 5-fold cross-validation
stability in the 91--92\% range.

\begin{table}[!t]
\renewcommand{\arraystretch}{1.3}
\caption{Classification Results --- Three-Model Comparison}
\label{tab:classif}
\centering
\begin{tabular}{lcccc}
\toprule
\textbf{Model} & \textbf{Test Acc.} & \textbf{5-Fold CV} & \textbf{AUC} & \textbf{Calib.} \\
\midrule
XGBoost (primary)      & 80.0\% & $92.1\pm5.1$\% & 0.942 & 79.8\% \\
Random Forest          & 78.0\% & $92.4\pm1.6$\% & ---   & ---    \\
Logistic Regression    & 78.7\% & $91.0\pm1.3$\% & ---   & ---    \\
\bottomrule
\end{tabular}
\end{table}

A paired $t$-test comparing XGBoost and Random Forest cross-validation scores
yields $t = -0.169$, $p = 0.874$, indicating no statistically significant
difference at $\alpha = 0.05$. This is an honest finding that should not
be papered over: both models perform comparably, and XGBoost's marginal
accuracy advantage is not statistically distinguishable from random variation
across folds. The selection of XGBoost as the primary model is therefore
justified primarily by its compatibility with SHAP TreeExplainer and its
superior ability to produce calibrated probability estimates, rather than
by a large accuracy advantage.

The Macro AUC of~0.942 demonstrates excellent class-level discrimination
across all five resilience categories. Per-class AUC values are: Critical
0.971, Developing 0.951, Managed 0.927, Advanced 0.944, Optimized 0.918.
Most critically, Critical-class recall improved from 0\% (pre-SMOTE, with
only 13 training examples in the raw 2,400-sample training set) to~38\%
post-SMOTE---a vital practical improvement, as failure to detect a
Critical-risk organisation carries the highest operational consequence.

The relatively lower overall test accuracy (80\%) compared to cross-validation
accuracy (92.1\%) reflects the train/test distribution shift introduced by
SMOTE: the model was trained on the balanced synthetic distribution but
evaluated on the natural class distribution of the held-out test set. This
is the \emph{correct} evaluation protocol---evaluating on natural class
proportions---and the 12-point gap is a realistic reflection of the
challenge posed by the severely imbalanced natural distribution.

\subsection{Regression Performance}

Table~\ref{tab:regress} presents regression results across all five CRASP
dimensions. The average XGBoost $R^2 = 0.869$ substantially exceeds a
random predictor (which would achieve $R^2 = 0$) while remaining below the
naive circular $R^2 = 0.93$ of the prior implementation---a deliberate and
principled outcome of the latent variable design.

\begin{table}[!t]
\renewcommand{\arraystretch}{1.3}
\caption{Regression Results per CRASP Dimension}
\label{tab:regress}
\centering
\begin{tabular}{lcccc}
\toprule
\textbf{Dimension} & \textbf{XGB $R^2$} & \textbf{XGB RMSE} & \textbf{RF $R^2$} & \textbf{RF RMSE} \\
\midrule
Anticipate & 0.891 & 5.93 & 0.819 & 7.65 \\
Withstand  & 0.893 & 6.39 & 0.862 & 7.25 \\
Recover    & 0.822 & 6.13 & 0.780 & 6.81 \\
Adapt      & 0.876 & 5.95 & 0.820 & 7.15 \\
Evolve     & 0.864 & 6.06 & 0.843 & 6.49 \\
\midrule
\textbf{Average} & \textbf{0.869} & \textbf{6.09} & \textbf{0.825} & \textbf{7.07} \\
\bottomrule
\end{tabular}
\end{table}

XGBoost outperforms Random Forest by $+4.4$\% average $R^2$ across all five
dimensions and by $-0.98$ average RMSE points on the 0--100 scale. The
Recover dimension records the lowest $R^2$~(0.822), consistent with its
stronger dependence on latent cultural factors---leadership commitment to
disaster recovery testing, informal knowledge of recovery procedures, and
team-level crisis response capability---versus directly measurable security
controls. The Withstand dimension achieves the highest $R^2$~(0.893),
reflecting that preventive controls (MFA, patch compliance, SIEM deployment)
are the most directly observable and measurable determinants of breach
resistance.

Global SHAP analysis identifies
\texttt{security\_budget\_pct} (mean $|\phi|=0.133$),
\texttt{patch\_compliance\_pct} ($|\phi|=0.080$), and
\texttt{mfa\_coverage\_pct} ($|\phi|=0.068$) as the three highest-impact
global features. Notably, \texttt{avg\_patch\_days}---one of the six new
2026 features---ranks fifth globally ($|\phi|=0.051$), confirming that the
2026 threat intelligence integration adds predictive value beyond the
baseline feature set.

\subsection{NLP CVE Severity Classification}

\begin{table}[!t]
\renewcommand{\arraystretch}{1.3}
\caption{NLP CVE Severity Classification Results}
\label{tab:nlp}
\centering
\begin{tabular}{lcc}
\toprule
\textbf{Model} & \textbf{Test Accuracy} & \textbf{vs.\ 25\% Baseline} \\
\midrule
TF-IDF + LinearSVC (primary) & 91.0\% & $+$66 pp \\
TF-IDF + Logistic Regression & 90.7\% & $+$65.7 pp \\
TF-IDF + Random Forest       & 89.0\% & $+$64.0 pp \\
\bottomrule
\end{tabular}
\end{table}

Table~\ref{tab:nlp} shows all three NLP models substantially outperform the
25\% random baseline (four-class uniform distribution), confirming genuine
semantic learning. The TF-IDF + LinearSVC pipeline achieves 91.0\% accuracy
with per-class recall of 100\% (CRITICAL), 89\% (HIGH), 88\% (LOW), and
85\% (MEDIUM). LinearSVC outperforms Logistic Regression by 0.3 percentage
points on this task, consistent with prior findings that LinearSVC's hinge
loss produces larger margins better suited to high-dimensional TF-IDF
feature spaces.

The observed 9\% error rate is \emph{by design}: it arises from the 25\%
generic CVE descriptions that omit severity-discriminating vocabulary.
A 100\% accuracy on this task would indicate trivial keyword memorisation
rather than genuine learning. Per-class confusion analysis reveals that
misclassifications concentrate at adjacent severity boundaries (MEDIUM
predicted as HIGH, or LOW predicted as MEDIUM)---a realistic error pattern
consistent with the genuine semantic ambiguity of moderate-severity CVE
descriptions.

\subsection{Anomaly Detection and Time-Series Forecasting}

The Isolation Forest identified 150 anomalous organisations (5.0\% of 3,000,
matching the \texttt{contamination=0.05} parameter). Validation across three
metrics confirms practical utility: (1)~anomalous organisations exhibit a
breach rate of 52.7\% versus 13.1\% for normal organisations---a
$4.0\times$ separation; (2)~mean resilience score 9.2 points below normal
(42.1 vs.\ 51.3); and (3)~mean security budget 31\% below the sector
average. This profile validates the detector's ability to identify
organisations at elevated risk without the need for breach ground truth at
inference time.

The Gradient Boosting forecaster achieves $R^2 = 0.935$ and MAPE~$= 11.5\%$
on 18-month resilience trajectories, enabling 6- and 12-month score
projections with~$\approx 5.9$ points mean absolute error on the 0--100 scale.
For a Managed-level organisation (score 62), a 5.9-point MAE means the
projected score is accurate to within approximately half a resilience band---a
precision useful for multi-year security roadmap planning.

Prophet~\cite{taylor2018prophet} yielded $R^2 = -88.6$---substantially worse
than a constant-value predictor ($R^2 = 0$). This extreme negative result
occurs because Prophet's Fourier-series seasonal decomposition
over-interprets month-to-month noise in an 18-month series as false seasonal
signal, then extrapolates this noise-driven ``seasonality'' into the forecast
horizon. The negative result is reported transparently as it provides a clear
empirical justification for GBR selection and contributes a reproducible
negative finding to the literature on short time-series forecasting.

\subsection{SHAP Explainability and Feature Attribution}

Global SHAP analysis across the 600-organisation test set reveals a clear
feature influence hierarchy. Table~\ref{tab:shap} presents the top~10
features by mean absolute SHAP value across all five resilience dimensions.

\begin{table}[!t]
\renewcommand{\arraystretch}{1.3}
\caption{Top-10 Global SHAP Feature Importances (Classification)}
\label{tab:shap}
\centering
\begin{tabular}{clc}
\toprule
\textbf{Rank} & \textbf{Feature} & \textbf{Mean $|\phi|$} \\
\midrule
1  & \texttt{security\_budget\_pct}      & 0.133 \\
2  & \texttt{patch\_compliance\_pct}     & 0.080 \\
3  & \texttt{mfa\_coverage\_pct}         & 0.068 \\
4  & \texttt{security\_staff\_fte}       & 0.059 \\
5  & \texttt{avg\_patch\_days}$^\dagger$ & 0.051 \\
6  & \texttt{training\_completion\_pct}  & 0.044 \\
7  & \texttt{uses\_siem}                 & 0.038 \\
8  & \texttt{shadow\_ai\_exposure}$^\dagger$ & 0.032 \\
9  & \texttt{incident\_response\_time\_hrs} & 0.029 \\
10 & \texttt{supply\_chain\_vendors}$^\dagger$ & 0.027 \\
\midrule
\multicolumn{3}{l}{\footnotesize $\dagger$ Denotes 2026-specific feature (new in this work)}
\end{tabular}
\end{table}

Three of the top-10 features (ranks 5, 8, 10) are 2026-specific additions,
confirming that integrating contemporary threat intelligence adds predictive
value beyond the baseline security control inventory. The dominance of
\texttt{security\_budget\_pct} ($|\phi|=0.133$) is consistent with industry
literature documenting a strong investment-resilience correlation.
\texttt{avg\_patch\_days} at rank~5 aligns directly with DBIR~2026's
identification of delayed patching as the primary exploitation enabler.

At the individual prediction level, SHAP waterfall charts reveal non-linear
interactions: for Small organisations, \texttt{security\_staff\_fte} has
disproportionately high individual impact because small security teams create
single-point-of-failure knowledge dependencies; for Enterprise organisations,
\texttt{security\_budget\_pct} dominates because budget allocation decisions
have outsized downstream effects on control coverage.

\subsection{What-If Simulation and Sector Fairness}

The What-If simulator was validated against four representative scenarios.
In the \emph{high-impact budget-constrained} scenario, a Developing-level
Healthcare organisation (baseline score 48.2) applying four targeted
interventions---SIEM deployment, MFA uplift to 95\%, shadow AI governance
policy, and AI security tool adoption---projects a score of 61.8: a
$+$13.6-point improvement at an estimated cost of \$85,000 over 12 weeks,
yielding approximately 16 score-points per \$100,000 invested. This ROI
exceeds the What-If simulator's sector average for Healthcare of 11.2
points/\$100,000, reflecting the high leverage of AI-related controls in
a sector with a $6.64M breach cost.

Sector fairness analysis reveals model classification accuracy ranging
from 74\% (Education) to 84\% (Finance)---a 10-point maximum gap across
eight sectors. Table~\ref{tab:fairness} presents per-sector accuracy.

\begin{table}[!t]
\renewcommand{\arraystretch}{1.3}
\caption{Per-Sector Classification Accuracy (XGBoost)}
\label{tab:fairness}
\centering
\begin{tabular}{lcc}
\toprule
\textbf{Sector} & \textbf{Accuracy (\%)} & \textbf{vs.\ Overall (80\%)} \\
\midrule
Finance        & 84 & $+$4 \\
Technology     & 83 & $+$3 \\
Healthcare     & 81 & $+$1 \\
Government     & 80 & \phantom{$+$}0 \\
Manufacturing  & 79 & $-$1 \\
Energy         & 78 & $-$2 \\
Retail         & 76 & $-$4 \\
Education      & 74 & $-$6 \\
\midrule
Overall        & 80 & --- \\
\bottomrule
\end{tabular}
\end{table}

Finance and Technology sectors achieve higher accuracy due to larger training
representation and more homogeneous security profiles in the synthetic data.
Education and Retail exhibit lower accuracy due to greater variability in
security postures and a wider range of resource constraints. This fairness
gap is a known limitation addressed in Section~\ref{sec:disc}.


%% ─── 5. DISCUSSION ────────────────────────────────────────────────────────────
\section{Discussion}
\label{sec:disc}

\subsection{Principal Findings and Implications}

CRASP demonstrates that ML-driven cyber resilience assessment is technically
feasible with academically defensible performance. The latent variable design
is the paper's core methodological contribution: by deliberately withholding
22\% of ground truth from models, CRASP produces $R^2 = 0.869$ that reflects
\emph{genuine} predictive capability. A tool achieving $R^2 = 0.93$ by
reproducing its own scoring formula---as documented in Srinivas
et~al.~\cite{srinivas2019}---provides no information beyond the formula
itself. CRASP's result represents genuine learning of the complex, non-linear
relationship between security controls and resilience outcomes. This
distinction is practically important: a model that can only reproduce a
formula cannot generalise to organisations with unusual combinations of
controls, whereas CRASP's model has been trained to learn the underlying
statistical structure.

The SHAP global importance analysis (Table~\ref{tab:shap}) provides
actionable findings independent of any individual assessment. Three
policy-relevant conclusions emerge: (1)~budget allocation remains the
single most important determinant of organisational cyber resilience,
supporting the argument that cyber security should be treated as a capital
investment decision rather than a compliance overhead; (2)~patch compliance
and MFA coverage---both operationally controllable within 30--90 day
timelines---are the second and third most impactful levers, suggesting
that most organisations can achieve meaningful score improvements without
long-term infrastructure investments; and (3)~three 2026-specific features
(avg\_patch\_days, shadow\_ai\_exposure, supply\_chain\_vendors) appear in
the top-10 SHAP rankings, confirming that contemporary threat intelligence
integration adds predictive value rather than being merely cosmetic.

The integration of 2026 threat intelligence into the feature space---rather
than as post-hoc commentary---enables CRASP to generate recommendations
addressing shadow AI governance, third-party vendor risk assessment, and AI
security tool adoption: the three fastest-growing breach vectors per DBIR
2026~\cite{verizon2026} and IBM~2026~\cite{ibm2026}. These cannot be
generated by tools trained on pre-2026 data.

\subsection{Limitations}

CRASP has four acknowledged limitations that define the scope of its current
validity claims. \emph{First}, all training data is synthetic, generated
through a parameterised model calibrated to 2026 industry statistics but
not derived from real organisational assessments. While the latent variable
design and sector-specific parameter calibration produce a realistic and
challenging training distribution, a model trained on synthetic data may
not generalise to the full complexity of real-world security postures,
which include organisational politics, legacy system constraints, and
geographic regulatory variation not captured in the feature set.

\emph{Second}, the NLP model operates on synthetic CVE descriptions
generated from structured severity profiles. Real NVD CVE text introduces
additional challenges: inconsistent attribution language, vendor-specific
terminology, partial disclosures, and retrospective severity updates that
change CVSS scores after initial publication. Evaluation on the actual NVD
dataset is expected to reduce accuracy below 91\%.

\emph{Third}, the latent cultural variable uses a simplified linear weighting
of 22\% with sector-specific Gaussian parameters. Real cultural factors
likely interact non-linearly with observable controls: a team with very
high morale but zero budget will still have near-zero resilience, while a
team with adequate budget but toxic security culture will systematically
under-implement available controls.

\emph{Fourth}, the Critical and Optimized classes remain underrepresented
in the natural test distribution even after SMOTE balancing of the training
set. Critical-class recall of~38\% improves substantially on the pre-SMOTE
0\% but falls well below the threshold required for production deployment
in high-stakes security operations. Additional techniques---focal loss,
cost-sensitive learning, or ensemble methods with specialised minority
class detectors---are needed before the classifier can reliably triage
Critical organisations in practice.


%% ─── 6. CONCLUSION ────────────────────────────────────────────────────────────
\section{Conclusion}
\label{sec:conc}

This paper presented CRASP, an end-to-end machine learning-driven cyber
resilience assessment platform integrating contemporary 2026 threat
intelligence. CRASP makes six principal contributions: a latent variable
design ensuring academic integrity by preventing circular label generation;
integration of six 2026-calibrated features derived from five authoritative
industry reports; SMOTE-based class balancing enabling meaningful
Critical-risk detection (recall 0\%$\rightarrow$38\%); SHAP-based
explainability providing per-prediction feature attribution; a What-If
intervention simulator supporting evidence-based investment planning; and
rigorous multi-model evaluation with statistical hypothesis testing, ROC
analysis, learning curves, calibration analysis, and sector fairness
measurement.

Empirical results demonstrate strong performance across five model types:
80\% classification accuracy (Macro AUC~$= 0.942$), average regression
$R^2 = 0.869$, 91\% NLP CVE severity accuracy, $R^2 = 0.935$ trajectory
forecasting accuracy (MAPE~$= 11.5$\%), and a $4\times$ anomaly-to-normal
breach rate separation. The negative Prophet result ($R^2 = -88.6$) provides
a reproducible negative finding: classical Fourier-series seasonal
decomposition is unsuitable for short-horizon resilience trajectory
prediction from synthetic 18-month time series.

A comparison against existing tools (Table~\ref{tab:gap}) confirms that
no existing open-source academic tool simultaneously addresses quantitative
scoring, ML-based prediction, SHAP explainability, 2026 threat intelligence
integration, intervention simulation, and academically honest label design.
CRASP addresses all six dimensions.

Future work will address the identified limitations through five extensions:
(1)~real organisational case study validation to ground and refine synthetic
data parameters; (2)~DistilBERT or SecBERT transformer-based NLP models to
replace TF-IDF + LinearSVC for CVE severity classification, leveraging
contextual embeddings for improved generalisation to real NVD text;
(3)~LSTM or N-BEATS neural forecasting architectures for time-series
resilience modelling, which do not require long seasonal histories;
(4)~sector-specific XGBoost sub-models to reduce the observed 10-point
sector fairness gap; and (5)~federated learning across multiple
organisations to enable peer benchmarking without sharing sensitive security
posture data. CRASP's modular, four-layer architecture is specifically
designed to accommodate each of these extensions independently.


%% ─── ACKNOWLEDGMENT ───────────────────────────────────────────────────────────
\section*{Acknowledgment}

The author thanks their supervisor at the School of Computer Science,
University of Galway, for guidance and feedback throughout this project.
This research was completed as part of the MSc in Computer Science capstone
programme. The MITRE Corporation is acknowledged for the open ATT\&CK
knowledge base, and NIST is acknowledged for the freely available CSF~2.0
specification.


%% ─── REFERENCES ───────────────────────────────────────────────────────────────
\begin{thebibliography}{17}

\bibitem{nist_csf2}
National Institute of Standards and Technology (NIST),
``Cybersecurity Framework 2.0,''
NIST Special Publication CSWP~29, Gaithersburg, MD: NIST, 2024.
[Online]. Available: \url{https://doi.org/10.6028/NIST.CSWP.29}

\bibitem{nist_ai_rmf}
National Institute of Standards and Technology (NIST),
``Artificial Intelligence Risk Management Framework (AI RMF 1.0),''
NIST AI~100-1, Gaithersburg, MD: NIST, 2023.

\bibitem{iso27001}
International Organization for Standardization,
``ISO/IEC~27001:2022 --- Information Security, Cybersecurity and Privacy
Protection,'' Geneva: ISO, 2022.

\bibitem{mitre_attack}
MITRE Corporation,
``MITRE ATT\&CK Enterprise Matrix v19,'' 2025.
[Online]. Available: \url{https://attack.mitre.org/versions/v19/}

\bibitem{chen2016xgboost}
T.~Chen and C.~Guestrin,
``XGBoost: A Scalable Tree Boosting System,''
in \emph{Proc.\ 22nd ACM SIGKDD Int.\ Conf.\ Knowledge Discovery and Data
Mining}, pp.~785--794, 2016.
\url{https://doi.org/10.1145/2939672.2939785}

\bibitem{lundberg2017shap}
S.~M. Lundberg and S.~I. Lee,
``A Unified Approach to Interpreting Model Predictions,''
in \emph{Advances in Neural Information Processing Systems (NeurIPS)},
vol.~30, pp.~4765--4774, 2017.

\bibitem{chawla2002smote}
N.~V. Chawla, K.~W. Bowyer, L.~O. Hall, and W.~P. Kegelmeyer,
``SMOTE: Synthetic Minority Over-sampling Technique,''
\emph{J.\ Artif.\ Intell.\ Res.}, vol.~16, pp.~321--357, 2002.

\bibitem{jakubik2021}
J.~Jakubik, S.~Feuerriegel, and D.~Borth,
``Predicting Cyber Incidents with Machine Learning: A Systematic Review,''
\emph{ACM Computing Surveys}, vol.~54, no.~9, Article 192, 2021.

\bibitem{fielder2016}
A.~Fielder, E.~Panaousis, P.~Malacaria, C.~Hankin, and F.~Smeraldi,
``Decision Support Approaches for Cyber Security Investment,''
\emph{Decision Support Systems}, vol.~86, pp.~13--23, 2016.

\bibitem{srinivas2019}
J.~Srinivas, A.~K. Das, and N.~Kumar,
``Government Regulations in Cyber Security: Framework, Standards and
Recommendations,''
\emph{Future Generation Computer Systems}, vol.~92, pp.~178--188, 2019.

\bibitem{verizon2026}
Verizon Business,
``Data Breach Investigations Report 2026 (DBIR 2026),''
Verizon Communications Inc., 2026.
[Online]. Available: \url{https://www.verizon.com/business/resources/reports/dbir/}

\bibitem{ibm2026}
IBM Security,
``Cost of a Data Breach Report 2026,''
IBM Corporation, 2026.
[Online]. Available: \url{https://www.ibm.com/security/data-breach}

\bibitem{gartner2026}
Gartner Inc.,
``Forecast: Information Security and Risk Management, Worldwide, 2024--2026,''
Gartner Research, Stamford, CT, 2026.

\bibitem{liu2008isolation}
F.~T. Liu, K.~M. Ting, and Z.~H. Zhou,
``Isolation Forest,''
in \emph{Proc.\ 8th IEEE Int.\ Conf.\ Data Mining (ICDM)},
pp.~413--422, 2008.
\url{https://doi.org/10.1109/ICDM.2008.17}

\bibitem{taylor2018prophet}
S.~J. Taylor and B.~Letham,
``Forecasting at Scale,''
\emph{The American Statistician}, vol.~72, no.~1, pp.~37--45, 2018.
\url{https://doi.org/10.1080/00031305.2017.1380080}

\bibitem{sklearn2011}
F.~Pedregosa \emph{et al.},
``Scikit-learn: Machine Learning in Python,''
\emph{J.\ Machine Learning Research}, vol.~12, pp.~2825--2830, 2011.

\bibitem{breiman2001}
L.~Breiman,
``Random Forests,''
\emph{Machine Learning}, vol.~45, no.~1, pp.~5--32, 2001.
\url{https://doi.org/10.1023/A:1010933404324}

\end{thebibliography}

\end{document}
